% ============================================================
\part{Die Idee der Robotik-Werkstatt}
% ============================================================

\chapter{Wir wollen einen Roboter programmieren -- aber womit eigentlich?}

\kapitelziel{
Nach diesem Kapitel verstehst du, warum Robotikentwicklung nicht nur aus
Programmcode besteht. Du lernst die wichtigsten Werkzeuge kennen, die später
immer wieder auftauchen: Linux, Terminal, Editor, Git, Docker,
Programmiersprachen und ROS2.
}
\kapitelstruktur


\section{Situation}

Stell dir vor, du möchtest einen mobilen Roboter entwickeln. Er soll Hindernisse
erkennen, Motoren ansteuern, Sensordaten auswerten und eine Karte seiner
Umgebung erstellen.

Die naheliegende Frage lautet:

\begin{quote}
Welche Programmiersprache brauche ich dafür?
\end{quote}

Diese Frage ist verständlich, aber zu eng. Ein Roboter entsteht nicht nur durch
eine Programmiersprache. Er entsteht durch ein Zusammenspiel vieler Werkzeuge.

\begin{aufgabeBox}
Wir wollen herausfinden, welche technische Grundausstattung wir brauchen,
um später Roboter mit ROS2 entwickeln zu können.
\end{aufgabeBox}

\section{Erster Lösungsversuch}

Man könnte sagen:

\begin{quote}
Ich installiere einfach Python und schreibe ein Programm.
\end{quote}

Für ein erstes kleines Programm reicht das. Für Robotik reicht es nicht.

Ein Roboterprojekt braucht meistens:
\begin{itemize}
  \item ein Betriebssystem,
  \item einen Editor,
  \item ein Terminal,
  \item eine Programmiersprache,
  \item eine Möglichkeit zur Versionsverwaltung,
  \item eine reproduzierbare Entwicklungsumgebung,
  \item Kommunikationsmechanismen zwischen Programmen,
  \item später ein Robotik-Framework wie ROS2.
\end{itemize}

\begin{problemBox}
Das eigentliche Problem ist nicht: \enquote{Wie schreibe ich eine einzelne Codezeile?}

Das eigentliche Problem ist: \enquote{Wie baue ich eine stabile Arbeitsumgebung,
in der Robotiksoftware zuverlässig entstehen kann?}
\end{problemBox}

\section{Was wir dafür lernen müssen}

\begin{wissenBox}
Für die Robotik-Werkstatt brauchen wir zunächst einen Überblick über:
\begin{itemize}
  \item Betriebssystem und Entwicklungsumgebung,
  \item Terminal und Quellcode,
  \item Versionsverwaltung (Git),
  \item Container (Docker),
  \item Robotik-Framework (ROS2).
\end{itemize}
\end{wissenBox}

\section{Die Werkzeuge im Überblick}

\begin{lstlisting}[style=text]
Robotik-Werkstatt
├── Computer
├── Linux
├── Terminal
├── VS Code
├── Git
├── Docker
├── Python
├── C++
└── später ROS2
\end{lstlisting}

\subsection{Linux}
Linux ist das Betriebssystem, auf dem sehr viele Robotiksysteme laufen. Es ist
stabil, gut automatisierbar und in der technischen Entwicklung weit verbreitet.

\subsection{Terminal}
Das Terminal ist die direkte Texteingabe für den Computer. Man startet Programme,
ruft Befehle auf, prüft Dateien und automatisiert Abläufe.

\subsection{VS Code}
Visual Studio Code ist Editor und Entwicklungsumgebung. Dort schreiben wir
Code, öffnen Projektordner, nutzen das Terminal und verwalten Dateien.

\subsection{Git}
Git merkt sich Entwicklungsstände. Wenn etwas kaputtgeht, kann man Änderungen
nachvollziehen oder zurückspringen.

\subsection{Docker}
Docker macht Entwicklungsumgebungen reproduzierbar. Was auf einem Rechner
läuft, soll auch auf einem anderen laufen.

\subsection{Python und C++}
Python ist einfach zu lernen und in ROS2 sehr gut nutzbar. C++ ist in der
professionellen Robotik wichtig für Geschwindigkeit und hardwarenahe Entwicklung.

\subsection{ROS2}
ROS2 ist kein Ersatz für Python oder C++. ROS2 ist ein Framework, mit dem
viele Robotikprogramme miteinander kommunizieren können.

\section{Konkrete Lösung}

\begin{loesungBox}
Für diese Lernreihe verwenden wir folgende Grundstruktur:
\begin{enumerate}
  \item Band 1: Linux, Docker und Entwicklungswerkzeuge.
  \item Band 2: Python für Robotik.
  \item Band 3: ROS2 mit Docker.
  \item Band 4: C++ für Robotik.
  \item Band 5: Mathematik für mobile Roboter.
  \item Band 6: Mathematik für Manipulatoren.
\end{enumerate}
\end{loesungBox}

\begin{merksatzBox}
Wer Roboter entwickeln will, braucht zuerst eine Werkstatt. Diese Werkstatt
besteht aus Betriebssystem, Terminal, Editor, Versionsverwaltung,
Containertechnik und später ROS2.
\end{merksatzBox}

\begin{fehlerBox}
\begin{itemize}
  \item \textbf{Fehler:} Direkt mit ROS2 starten, ohne Linux und Terminal zu verstehen.\\
        \textbf{Folge:} Schon einfache Fehlermeldungen werden unverständlich.
  \item \textbf{Fehler:} Alles lokal installieren, ohne Dokumentation.\\
        \textbf{Folge:} Die Umgebung lässt sich kaum reproduzieren.
  \item \textbf{Fehler:} Kein Git verwenden.\\
        \textbf{Folge:} Änderungen sind nicht nachvollziehbar.
\end{itemize}
\end{fehlerBox}

\begin{uebungBox}
\textbf{Nachmachen:} Schreibe eine Liste der Werkzeuge, die du bereits installiert hast.

\textbf{Verändern:} Ergänze zu jedem Werkzeug: \enquote{Dieses Werkzeug brauche ich, weil\ldots}

\textbf{Anwenden:} Beschreibe mit eigenen Worten, warum ein Roboterprojekt
nicht nur aus einem Python-Programm besteht.
\end{uebungBox}

\begin{quizBox}
\begin{enumerate}
  \item Ist ROS2 eine Programmiersprache?
  \item Wofür braucht man Git?
  \item Warum ist Docker in der Robotik nützlich?
  \item Warum reicht es nicht, nur Python zu lernen?
\end{enumerate}
\textbf{Lösungen:}
\begin{enumerate}
  \item Nein. ROS2 ist ein Robotik-Framework (Middleware).
  \item Zur Versionsverwaltung und Nachvollziehbarkeit von Änderungen.
  \item Für reproduzierbare Entwicklungsumgebungen.
  \item Weil Robotik zusätzlich Betriebssysteme, Schnittstellen und Hardware erfordert.
\end{enumerate}
\end{quizBox}

\begin{haubeBox}
\subsection*{Unter der Haube: Robotikentwicklung als Schichtenmodell}

\begin{lstlisting}[style=text]
Anwendungsebene
    ↓
Robotik-Framework (ROS2)
    ↓
Programmiersprachen (Python / C++)
    ↓
Bibliotheken und Laufzeitumgebungen
    ↓
Betriebssystem (Linux)
    ↓
Treiber und Schnittstellen
    ↓
Hardware
\end{lstlisting}

Ein professioneller Entwickler fragt nicht nur: \enquote{Warum geht mein Code nicht?}
Er fragt: \enquote{Auf welcher Ebene liegt das Problem?} -- denn der Fehler kann
auf jeder dieser Schichten liegen.
\end{haubeBox}

\begin{robotikBox}
Später in ROS2 werden einzelne Roboterfunktionen auf viele Programme verteilt.
Eine Kamera-Node verarbeitet Bilder. Eine Lidar-Node liefert Abstandsdaten.
Eine Motor-Node steuert Antriebe.

Ohne Grundverständnis für Prozesse, Dateien, Netzwerke und Entwicklungsumgebungen
wird ROS2 schnell unübersichtlich.
\end{robotikBox}

% ------------------------------------------------------------

\chapter{Warum verwenden so viele Roboter Linux?}

\kapitelziel{
Nach diesem Kapitel verstehst du, warum Linux in der Robotik so verbreitet
ist und warum wir uns in Band 1 nicht auf Windows als Hauptumgebung stützen.
}
\kapitelstruktur


\section{Situation}

Du suchst im Internet nach ROS2, Raspberry Pi, Docker oder mobiler Robotik.
Sehr schnell fällt auf: Viele Anleitungen verwenden Ubuntu oder eine andere
Linux-Distribution.

\begin{aufgabeBox}
Wir wollen verstehen, warum Linux für Robotikentwicklung so wichtig ist.
\end{aufgabeBox}

\begin{problemBox}
Wir müssen einordnen, ob Linux nur eine Vorliebe von Entwicklern ist
oder ob es technische Gründe dafür gibt.
\end{problemBox}

\begin{wissenBox}
Wir brauchen ein Grundverständnis für: Betriebssysteme, Open Source,
Automatisierbarkeit, Embedded-Systeme, Server- und Netzwerktechnik,
ROS2-Unterstützung.
\end{wissenBox}

\section{Konkrete Gründe für Linux in der Robotik}

\begin{itemize}
  \item Linux läuft auf kleinen Rechnern wie Raspberry Pi.
  \item Linux läuft auf Industrie-PCs und Steuerungsrechnern.
  \item Linux lässt sich gut automatisieren und skripten.
  \item Linux bietet starken Netzwerkkommunikations-Stack.
  \item Viele Robotikbibliotheken sind für Linux zuerst verfügbar.
  \item ROS2 wird unter Ubuntu offiziell unterstützt.
  \item Hardware-Geräte erscheinen als Dateien (\texttt{/dev/ttyUSB0}).
\end{itemize}

\begin{loesungBox}
Für unsere Lernreihe gilt: Wir verwenden Linux, weil es der praktischste
Einstieg in professionelle ROS2-Entwicklung ist. Windows kann als Hostsystem
dienen -- über Docker Desktop oder WSL läuft dann eine Linux-Umgebung darin.
\end{loesungBox}

\begin{merksatzBox}
Linux ist in der Robotik nicht deshalb wichtig, weil es exotisch ist, sondern
weil es stabil, automatisierbar und technisch gut kontrollierbar ist.
\end{merksatzBox}

\begin{haubeBox}
\subsection*{Unter der Haube: Linux als Schnittstelle zwischen Software und Hardware}

Ein Betriebssystem verwaltet Ressourcen: Prozessorzeit, Arbeitsspeicher,
Dateien, Benutzerrechte, Netzwerk, Hardwaregeräte.

Für Robotik ist besonders wichtig, dass Linux viele Geräte als Dateien
sichtbar macht. Ein serieller USB-Adapter erscheint als \texttt{/dev/ttyUSB0}.
Ein Python-Programm kann dann über diese Gerätedatei mit einem Mikrocontroller
kommunizieren -- ohne Treiber-DLLs, ohne COM-Port-Konfiguration.
\end{haubeBox}

\begin{robotikBox}
Wenn später ein ROS2-Programm mit einem Lidar, einer Kamera oder einem
Mikrocontroller spricht, läuft darunter oft ein Linux-Gerätetreiber. ROS2
ist dann nur die obere Kommunikationsschicht. Darunter steht weiterhin das
Betriebssystem.
\end{robotikBox}
